% Part: model-theory
% Chapter: models-of-arithmetic
% Section: non-standard-models

\documentclass[../../../include/open-logic-section]{subfiles}

\begin{document}

\olfileid[pt]{mod}{mar}{mdq}
\section{Modelos de $\Th{Q}$}

\begin{explain}
Sabemos que existem !!{structure}s não padrão que tornam verdadeiras
as mesmas !!{sentence}s que~$\Struct{N}$, isto é, que são modelos
de~$\Th{TA}$. Como $\Sat{N}{\Th{Q}}$, todo modelo de~$\Th{TA}$ também
é um modelo de~$\Th{Q}$. $\Th{Q}$ é muito mais fraca que~$\Th{TA}$;
por exemplo, $\Th{Q} \Proves/ \lforall[x][\lforall[y][\eq[(x +
      y)][(y+x)]]]$. Teorias mais fracas são mais fáceis de satisfazer:
elas têm mais modelos. Por exemplo, $\Th{Q}$ tem modelos que tornam
$\lforall[x][\lforall[y][\eq[(x + y)][(y+x)]]]$ falsa, mas estes não
podem ser também modelos de~$\Th{TA}$, nem de $\Th{PA}$. Os modelos de
$\Th{Q}$ também são relativamente simples: podemos especificá-los
explicitamente.
\end{explain}

\begin{ex}
\ollabel{ex:model-K-of-Q}
Considere a !!{structure}~$\Struct{K}$ com domínio $\Domain{K} = \Nat
\cup \{a\}$ e interpretações
\begin{align*}
  \Assign{\Obj{0}}{K} & = 0\\
  \Assign{\prime}{K}(x) & =
  \begin{cases}
    x+1 & \text{se $x\in \Nat$}\\
    a & \text{se $x = a$}
  \end{cases}\\
  \Assign{+}{K}(x, y) & =
  \begin{cases}
    x+y & \text{se $x$, $y \in\Nat$}\\
    a & \text{caso contrário}
  \end{cases}\\
  \Assign{\times}{K}(x, y) & =
  \begin{cases}
    xy & \text{se $x$, $y \in\Nat$}\\
    0 & \text{se $x = 0$ ou $y = 0$}\\
    a & \text{caso contrário}\\
  \end{cases}\\
  \Assign{<}{K} & =
  \Setabs{\tuple{x,y}}{x, y \in \Nat \text{ e } x<y} \cup
  \Setabs{\tuple{x,a}}{x \in \Domain{K}}
\end{align*}
Para mostrar que $\Sat{K}{\Th{Q}}$, precisamos verificar que todos os
axiomas de~$\Th{Q}$ são verdadeiros em~$\Struct{K}$. Por conveniência,
escrevamos $x^\nssucc$ para $\Assign{\prime}{K}(x)$ (o ``sucessor'' de
$x$ em~$\Struct{K}$), $x \nsplus y$ para $\Assign{+}{K}(x, y)$ (a
``soma'' de $x$ e $y$ em~$\Struct{K}$), $x \nstimes y$ para
$\Assign{\times}{K}(x, y)$ (o ``produto'' de $x$ e~$y$ em
$\Struct{K}$) e $x \nsless y$ para $\tuple{x,y} \in \Assign{<}{K}$.
Com essas abreviações, podemos apresentar as operações em~$\Struct{K}$
de modo mais claro como
\[
\begin{array}{c|c}
  x & x^\nssucc \\
  \hline
  n & n+1 \\
  a & a
\end{array}
\qquad
\begin{array}{c|ccc}
  x \nsplus y & 0 & m & a \\
  \hline
  0 & 0 & m & a \\
  n & n & n+m & a \\
  a & a & a & a \\
\end{array}
\qquad
\begin{array}{c|ccc}
  x \nstimes y & 0 & m & a \\
  \hline
  0 & 0 & 0 & 0 \\
  n & 0 & nm & a \\
  a & 0 & a & a \\
\end{array}
\]
Temos $n \nsless m$ se, e somente se, $n<m$ para $n$, $m \in \Nat$, e
$x \nsless a$ para todo~$x \in \Domain{K}$.

$\Sat{K}{\lforall[x][\lforall[y][(\eq[x'][y'] \lif \eq[x][y])]]}$,
pois $\nssucc$ é !!{injective}. $\Sat{K}{\lforall[x][\eq/[\Obj
0][x']]}$, pois $0$ não é sucessor segundo $\nssucc$ em~$\Struct{K}$.
$\Sat{K}{\lforall[x][(\eq[x][\Obj 0] \lor \lexists[y][\eq[x][y']])]}$,
pois, para todo $n>0$, $n = (n-1)^\nssucc$, e $a = a^\nssucc$.

$\Sat{K}{\lforall[x][\eq[(x + \Obj 0)][x]]}$, pois $n \nsplus 0 = n+0
= n$, e $a\nsplus 0 = a$ pela definição de~$\nsplus$.
$\Sat{K}{\lforall[x][\lforall[y][\eq[(x + y')][(x + y)']]]}$ é um
pouco mais trabalhoso. Se $n$ e $m$ são ambos padrão, temos:
\begin{align*}
(n \nsplus m^\nssucc) & = (n+(m+1)) = (n+m)+1 = (n \nsplus m)^\nssucc  
\intertext{pois $\nsplus$ e $^\nssucc$ coincidem com $+$ e $\prime$ nos
  números padrão. Suponha agora que $x \in \Domain{K}$. Então}
(x \nsplus a^\nssucc) & = (x \nsplus a) = a = a^\nssucc = (x \nsplus a)^\nssucc
\intertext{O caso restante ocorre quando $y \in \Domain{K}$, mas $x =
  a$. Aqui também devemos distinguir os casos conforme $y = n$ seja
  padrão ou $y = b$:}
(a \nsplus n^\nssucc) & = (a \nsplus (n+1)) = a = a^\nssucc = (a \nsplus n)^\nssucc\\
(a \nsplus a^\nssucc) & = (a \nsplus a) = a = a^\nssucc = (a \nsplus a)^\nssucc
\end{align*}
Isto é, naturalmente, um pouco mais detalhado do que o necessário. Por
exemplo, como $a \nsplus z = a$ seja qual for $z$, podemos concluir de
imediato que $a \nsplus a^\nssucc = a$. Os axiomas restantes podem ser
verificados do mesmo modo.

Assim, $\Struct{K}$ é um modelo de~$\Th{Q}$. Sua ``adição''~$\nsplus$
também é comutativa. Mas há outras sentenças verdadeiras em~$\Struct{N}$
e falsas em~$\Struct{K}$, e vice-versa. Por exemplo, $a \nsless a$,
logo $\Sat{K}{\lexists[x][x < x]}$ e
$\Sat/{K}{\lforall[x][\lnot x<x]}$. Isso mostra que $\Th{Q} \Proves/
\lforall[x][\lnot x < x]$.
\end{ex}

\begin{prob}
Demonstre que $\Struct{K}$ de \olref[mod][mar][mdq]{ex:model-K-of-Q}
satisfaz os axiomas restantes de~$\Th{Q}$,
\begin{align*}
  & \lforall[x][\eq[(x \times \Obj 0)][\Obj 0]] \tag{$!Q_6$}\\
  & \lforall[x][\lforall[y][\eq[(x \times y')][((x \times y) + x)]]] \tag{$!Q_7$}\\
  & \lforall[x][\lforall[y][(x < y \liff \lexists[z][\eq[(z' + x)][y]])]] \tag{$!Q_8$}
\end{align*}
Encontre !!a{sentence} que envolva apenas~$\prime$, seja verdadeira
em~$\Struct{N}$ e falsa em~$\Struct{K}$.
\end{prob}

\begin{ex}
\ollabel{ex:model-L-of-Q} Considere a !!{structure}~$\Struct{L}$ com
domínio $\Domain{L} = \Nat \cup \{a, b\}$ e interpretações
$\Assign{\prime}{L} = \nssucc$, $\Assign{+}{L} = \nsplus$ dadas por
\[
\begin{array}{c|c}
  x & x^\nssucc \\
  \hline
  n & n+1 \\
  a & a\\
  b & b
\end{array}
\qquad
\begin{array}{c|ccc}
  x \nsplus y & m & a & b\\
  \hline
  n & n+m & b & a\\
  a & a & b & a\\
  b & b & b & a
\end{array}
\]
Como $\nssucc$ é !!{injective}, $0$ não pertence à sua imagem, e todo
$x \in \Domain{L}$ diferente de~$0$ pertence, os axiomas
$!Q_1$--$!Q_3$ são verdadeiros em~$\Struct{L}$. Para todo $x$,
$x \nsplus 0 = x$, portanto $!Q_4$ também é verdadeiro. Para $!Q_5$,
considere $x \nsplus y^\nssucc$ e $(x \nsplus y)^\nssucc$. Eles são
iguais se $x$ e $y$ são ambos padrão, pois nesse caso $\nssucc$ e
$\nsplus$ coincidem com $\prime$ e $+$. Se $x$ é não padrão e $y$ é
padrão, temos $x \nsplus y^\nssucc = x = x^\nssucc = (x \nsplus
y)^\nssucc$. Se $x$ e $y$ são ambos não padrão, temos quatro casos:
\begin{align*}
&  a \nsplus a^\nssucc  = b = b^\nssucc = (a \nsplus a)^\nssucc\\
&  b \nsplus b^\nssucc  = a = a^\nssucc = (b \nsplus b)^\nssucc\\
&  b \nsplus a^\nssucc  = b = b^\nssucc = (b \nsplus y)^\nssucc\\
&  a \nsplus b^\nssucc  = a = a^\nssucc = (a \nsplus b)^\nssucc\\
\intertext{Se $x$ é padrão, mas $y$ é não padrão, temos}
&  n \nsplus a^\nssucc  = n \nsplus a = b = b^\nssucc = (n \nsplus a)^\nssucc\\
&  n \nsplus b^\nssucc  = n \nsplus b = a = a^\nssucc = (n \nsplus b)^\nssucc
\end{align*}
Logo, $\Sat{L}{!Q_5}$. No entanto, $a \nsplus 0 \neq 0 \nsplus a$, portanto
$\Sat/{L}{\lforall[x][\lforall[y][\eq[(x+y)][(y+x)]]]}$.
\end{ex}

\begin{prob}
Expanda $\Struct{L}$ de \olref[mod][mar][mdq]{ex:model-L-of-Q} para
incluir $\nstimes$ e $\nsless$, que interpretam~$\times$ e $<$. Mostre
que sua estrutura satisfaz os axiomas restantes de~$\Th{Q}$,
\begin{align*}
& \lforall[x][\eq[(x \times \Obj 0)][\Obj 0]] \tag{$!Q_6$}\\ &
  \lforall[x][\lforall[y][\eq[(x \times y')][((x \times y) + x)]]]
  \tag{$!Q_7$}\\ & \lforall[x][\lforall[y][(x < y \liff \lexists[z][\eq[(z'+x)][y]])]] \tag{$!Q_8$}
\end{align*}
\end{prob}

\begin{prob}
Em $\Struct{L}$ de \olref[mod][mar][mdq]{ex:model-L-of-Q}, $a^\nssucc
= a$ e $b^\nssucc = b$. Existe um modelo de~$\Th{Q}$ no qual
$a^\nssucc = b$ e $b^\nssucc = a$?
\end{prob}

\begin{explain}
Construímos explicitamente modelos de~$\Th{Q}$ nos quais os
!!{element}s não padrão ficam ``além'' dos elementos padrão. Na
verdade, isso é exigido pelos axiomas. Um !!{element} não padrão~$x$
não pode satisfazer ${} \nsless 0$, pois $\Th{Q} \Proves
\lforall[x][\lnot x<0]$ (veja \olref[inc][req][min]{lem:less-zero}).
Além disso, para todo $n$, $\Th{Q} \Proves \lforall[x][(x <
\num{n}' \lif (\eq[x][\num{0}] \lor \eq[x][\num{1}] \lor \dots \lor
\eq[x][\num{n}]))]$ (\olref[inc][req][min]{lem:less-nsucc}), portanto
não podemos ter $a \nsless n$ para nenhum~$n>0$.
\end{explain}

\end{document}
